\documentclass[11pt]{article}
\usepackage[a4paper,margin=28mm]{geometry}
\usepackage[T1]{fontenc}
\usepackage{lmodern,amsmath,amssymb,amsthm,mathtools,microtype}
\usepackage[colorlinks=true,linkcolor=blue,citecolor=blue,urlcolor=blue]{hyperref}
\newtheorem{theorem}{Theorem}
\newtheorem{lemma}{Lemma}
\newtheorem{remark}{Remark}
\newcommand{\R}{\mathbb R}
\newcommand{\Lop}{\mathcal L}
\newcommand{\norm}[1]{\lVert#1\rVert_\infty}
\title{Failure of comparison for a relaxed infinity fractional evolution}
\author{Kaiwen Jin\\\small\texttt{yui.ui.siki@gmail.com}}
\date{}
\begin{document}
\maketitle
\begin{abstract}
We consider the infinity fractional evolution in dimension $n\ge2$, with
$1/2<s<1$. At a zero spatial test gradient, the viscosity subsolution condition
uses the supremum of symmetric directional integrals, and the supersolution
condition uses their infimum. For this convention, comparison fails even for
radial subsolutions and supersolutions with the same compactly
supported initial data, attained uniformly. The construction uses two thin
annuli: a line tangent to the inner annulus collects substantially more mass
than a radial line. This permits a locally flat subsolution to rise faster
than a radial solution of the one-dimensional fractional heat equation.
\end{abstract}

\section{The comparison question}

Comparison for an evolution equation requires ordered initial data to keep
subsolutions below supersolutions. For a direction selected by the gradient,
its formulation at critical points matters. We consider the following precise
choice. Fix $n\ge2$, $1/2<s<1$, $C_s>0$ and $T>0$, and write $p=2s$.
For a unit vector $e$, define
\begin{equation}\label{eq:direction}
 D_e^s H(x)=C_s\int_0^\infty
 \frac{H(x+ze)+H(x-ze)-2H(x)}{z^{1+p}}\,dz.
\end{equation}
The integral is absolutely convergent for bounded $C^2$ functions with bounded
Hessian. The equation on $\R^n\times(0,T)$ is
$u_t=\Delta_\infty^s u$, with the direction $e=\nabla u/|\nabla u|$ where
$\nabla u\ne0$.

A viscosity subsolution is tested from above by a local smooth function
$\phi$; at the contact time, replace $\phi$ outside the spatial contact
neighborhood by the tested function. Denote this patch by $\widetilde\phi$.
The required inequalities at contact are
\begin{equation}\label{eq:weak}
 \phi_t\le
 \begin{cases}
 D_{\nabla\phi/|\nabla\phi|}^s\widetilde\phi,&\nabla\phi\ne0,\\
 \displaystyle\sup_{|e|=1}D_e^s\widetilde\phi,&\nabla\phi=0.
 \end{cases}
\end{equation}
For a supersolution use lower tests, reverse the inequality and use the
infimum at zero gradient. Suprema and infima are pointwise over all unit
directions. Initial data are bounded and uniformly continuous, and their
attainment is in the uniform norm on $\R^n$.

The infinity fractional Laplacian arose from nonlocal tug-of-war games in
Bjorland, Caffarelli and Figalli~\cite{BCF}. Their Section~2.3 distinguishes
weak symmetric directional conditions from the stronger zero-gradient
operator in their Definition~1.1, and discusses failure of elliptic comparison
for the weak conditions. For the evolution, del Teso, Endal, Jakobsen and
V\'azquez~\cite{DEJV} construct viscosity solutions; their Definition~4.3 and
Remark~4.4 use the relaxed envelopes in \eqref{eq:weak}. Their Theorem~2.6
compares constructed viscosity solutions with classical subsolutions and
supersolutions, while Section~9 leaves general comparison and uniqueness
open. The theorem below concerns exactly \eqref{eq:weak} with local patches.
It does not address the stronger zero-gradient definition or comparison
within a separately selected class of game limits.

\begin{theorem}\label{thm:main}
For every integer $n\ge2$, every $s\in(1/2,1)$, every $C_s>0$ and every
$T>0$, there are bounded functions $U,V$ on $\R^n\times[0,T]$, with
continuous bounded spatial derivatives through order two and continuous
bounded time derivative, such that $U$ is a viscosity subsolution and $V$ is
a viscosity supersolution in the convention \eqref{eq:weak}. They have the
same initial data $F\in C_c^\infty(\R^n)$, attained uniformly, but
$U(x_*,t_*)>V(x_*,t_*)$ for some $x_*\in\R^n$ and $0<t_*<T$.
In fact $V$ satisfies both viscosity inequalities. The functions are radial;
no monotonicity in the radius is imposed.
\end{theorem}

Thus the unrestricted subsolution--supersolution comparison implication is
false. The result does not construct two distinct full solutions. Its
mechanism is geometric: thin annuli have small mass along a radial line,
whereas a nearly tangent line stays close to one annulus for a much longer
segment. The supremum in \eqref{eq:weak} can use that segment on an open
region where the subsolution is spatially flat. Section~2 gives the estimates;
Section~3 constructs the global pair and verifies the actual local tests.
All constants below may depend on the fixed parameters. We use Lebesgue
measure, and $\norm{D^2H}$ denotes the uniform Euclidean operator norm.

\section{Two thin annuli}

For a real profile on the line set
\[
 \Lop h(r)=C_s\int_0^\infty
 \frac{h(r+z)+h(r-z)-2h(r)}{z^{1+p}}\,dz.
\]
Fix
\[
 \varepsilon=10^{-12},\qquad c_0=e^{-4/3},\qquad
 \psi(y)=\begin{cases}e^{-1/(1-y^2)},&|y|<1,\\0,&|y|\ge1,\end{cases}
\]
and define
\begin{equation}\label{eq:profiles}
 G(r)=\psi\bigl((|r|-1)/\varepsilon\bigr)
      +\psi\bigl((|r|-3)/\varepsilon\bigr),\qquad
 F(x)=G(|x|),\qquad b=\Lop G.
\end{equation}
The bump is smooth and flat at its support endpoints, $0\le\psi\le1$, and
$\psi\ge c_0$ on $[-1/2,1/2]$. Both $G$ and $F$ vanish near the origin
and are compactly supported smooth functions. Put
\[
 I=(1+\varepsilon,3-\varepsilon),\qquad
 a=\frac{C_s c_0\sqrt\varepsilon}{64}.
\]
Here $\varepsilon$ is a fixed part of the example, chosen small enough to
separate the radial and tangent estimates.

\begin{lemma}\label{lem:annuli}
The function $b$ is even and smooth, with every derivative bounded. Moreover,
\begin{equation}\label{eq:radialbounds}
 0<b(2)<a,\qquad b(1+\varepsilon)>2a,\qquad b(3-\varepsilon)>2a.
\end{equation}
For every $x$ with $|x|\in I$ there is a unit direction $e$ such that
\begin{equation}\label{eq:tangentbound}
 D_e^sF(x)\ge2a.
\end{equation}
\end{lemma}
\begin{proof}
For any bounded $C^2$ function $H$ with bounded Hessian, its symmetric
second difference is bounded by $\norm{D^2H}z^2$ and by $4\norm H$.
Splitting the jump integral at $z=1$ gives
\begin{equation}\label{eq:universalbound}
 |D_e^sH(x)|\le C_s\left(\frac{\norm{D^2H}}{2-p}
                         +\frac{4\norm H}{p}\right).
\end{equation}
Applying the same bounds to every derivative of $G$ justifies differentiation
under the integral by dominated convergence, and gives the stated properties
of $b$. For $r\in\overline I$, evenness and $G(r)=0$ yield
\begin{equation}\label{eq:bkernel}
 b(r)=C_s\int_0^\infty G(\rho)
 \left(|r-\rho|^{-1-p}+(r+\rho)^{-1-p}\right)d\rho.
\end{equation}
This remains finite at the endpoints by the preceding second-difference
bound. At $r=2$, the kernel on either bump is less than
$(4/3)^3+(4/11)^2<3$. Each bump has mass at most $2\varepsilon$, so
\[
 0<b(2)<12C_s\varepsilon<a.
\]
The final inequality follows from $768\sqrt\varepsilon<c_0$; for example,
$c_0>1/9>768\cdot10^{-6}$.
At the left endpoint restrict \eqref{eq:bkernel} to
$\rho\in[1,1+\varepsilon/2]$, and at the right endpoint to
$\rho\in[3-\varepsilon/2,3]$. There $G\ge c_0$, the relevant distance is
at most $\varepsilon$, and the interval length is $\varepsilon/2$. Hence
\[
 b(1+\varepsilon),\ b(3-\varepsilon)
 \ge \frac{C_s c_0}{2\varepsilon^p}>2a.
\]

To prove \eqref{eq:tangentbound}, write $r=|x|\in I$ and choose a unit
vector $\omega$ perpendicular to $x$, possible since $n\ge2$. Set
\[
 e=\sqrt{1-r^{-2}}\,\frac{x}{r}+r^{-1}\omega,
 \qquad z_0=\sqrt{r^2-1}.
\]
Then $|e|=1$, $x\cdot e=z_0$, and
\[
 |x-ze|^2=1+(z-z_0)^2.
\]
For $|z-z_0|\le\sqrt\varepsilon$, this radius belongs to
$[1,1+\varepsilon/2]$, so $F(x-ze)\ge c_0$. Also
$z_0>\sqrt\varepsilon$ and $z_0+\sqrt\varepsilon<4$.
Since $F(x)=0$ and $F\ge0$, integrating over this positive jump interval gives
\[
 D_e^sF(x)\ge 2C_s c_0\sqrt\varepsilon\,4^{-1-p}
 \ge\frac{C_s c_0\sqrt\varepsilon}{32}=2a.
\]
\end{proof}

\section{A global subsolution and a radial supersolution}

We first raise the low values of $b$ to a constant while leaving its large
values available to control the radial inequality. Choose a smooth
$\eta:\R\to[0,1]$ equal to $1$ on $(-\infty,0]$ and $0$ on $[1,\infty)$.
For instance, with $h(z)=e^{-1/z}$ for $z>0$ and $h(z)=0$ otherwise, take
$\eta(z)=h(1-z)/(h(z)+h(1-z))$. For $r\ge0$ set
\begin{equation}\label{eq:cutoff}
 \chi(r)=\begin{cases}\eta((b(r)-a)/a),&r\in I,\\0,&r\notin I,
 \end{cases}
 \qquad q(r)=b(r)+\chi(r)(a-b(r)),
\end{equation}
and extend $q$ evenly. By \eqref{eq:radialbounds}, $\chi$ is zero in collars
of both endpoints of $I$. Thus $q$ is smooth, bounded, and has bounded
derivatives of every order. Its strict upward modification has a useful form:
\begin{equation}\label{eq:flat}
 q(r)>b(r)\quad\Longrightarrow\quad
 |r|\in I,\quad G=0\ \hbox{and}\ q=a
 \ \hbox{on a neighborhood of }r.
\end{equation}
Indeed $q>b$ implies $b<a$, and then $\eta((b-a)/a)=1$ locally. Elsewhere
$q\le b$. In particular $q(2)=a$.

Set $Q(x)=q(|x|)$. To include the origin, note that an even $C_b^2$ profile
has $q'(0)=0$ and
\[
 \frac{q'(r)}r=\int_0^1q''(\theta r)\,d\theta\longrightarrow q''(0).
\]
The radial Hessian has eigenvalues $q''(r)$ and $q'(r)/r$, so it extends
continuously at zero as $q''(0)\operatorname{Id}$. Consequently $Q\in C_b^2$.
Define
\begin{equation}\label{eq:BM}
 B=C_s\left(\frac{\norm{D^2Q}}{2-p}+\frac{4\norm Q}{p}\right),
 \qquad M=B+1.
\end{equation}
Equation~\eqref{eq:universalbound} gives $|D_e^sQ|\le B$ in every direction.

\begin{lemma}\label{lem:sub}
The function $U(x,t)=F(x)+tQ(x)-Mt^2$ satisfies the pointwise subsolution
inequalities in \eqref{eq:weak} at every $x$ and every $t\ge0$.
\end{lemma}
\begin{proof}
If $x\ne0$, put $r=|x|$ and $e_r=x/r$. A nonzero gradient of $U$ is
parallel to $e_r$. For an even profile, $|x\pm ze_r|=|r\pm z|$ gives
\begin{equation}\label{eq:radialU}
 D_{e_r}^sU(x,t)=b(r)+t\Lop q(r)\ge b(r)-Bt.
\end{equation}
At the origin the same identity holds in every direction. Where $q\le b$,
\[
 U_t=q-2Mt\le b-Bt\le D_{e_r}^sU.
\]
This is the selected inequality at nonzero gradient, since $D_{-e_r}^s=D_{e_r}^s$.
At zero gradient its supremum is at least this radial value.
Where $q>b$, \eqref{eq:flat} makes $U$ locally spatially constant.
The tangent direction of Lemma~\ref{lem:annuli} gives
\[
 \sup_{|e|=1}D_e^sU\ge2a-Bt\ge a-2Mt=U_t.
\]
These two cases include every interface and the entire unbounded exterior.
\end{proof}

The supersolution comes from a linear evolution on the line, then a radial
lift. This lift will be used in the relaxed sense; radial monotonicity is
unnecessary for its radial directional identity.

\begin{lemma}\label{lem:heat}
There is a real even bounded function $w(r,t)$ on $\R\times[0,\infty)$,
with bounded continuous derivatives of every finite spatial and time order,
that satisfies $w_t=\Lop w$, $w(r,0)=G(r)$, and
$\sup_r|w(r,t)-G(r)|\to0$ as $t\downarrow0$.
For $V(x,t)=w(|x|,t)$, both pointwise relaxed inequalities hold.
\end{lemma}
\begin{proof}
With the convention $\widehat G(\xi)=\int_\R e^{-ir\xi}G(r)\,dr$, put
\begin{equation}\label{eq:heat}
 \kappa=2C_s\int_0^\infty\frac{1-\cos z}{z^{1+p}}\,dz,\qquad
 w(r,t)=\frac1{2\pi}\int_\R
 e^{ir\xi-\kappa t|\xi|^p}\widehat G(\xi)\,d\xi.
\end{equation}
The constant $\kappa$ is positive and finite, since $0<p<2$.
Integration by parts in the transform of the compactly supported smooth $G$
shows that $\int_\R(1+|\xi|)^k|\widehat G(\xi)|\,d\xi<\infty$ for every
$k\ge0$. Thus all differentiations in \eqref{eq:heat} have integrable
majorants independent of $r$ and $t\ge0$. Time derivatives at zero are right
derivatives. Fourier inversion for Schwartz functions, with the displayed
normalization, gives $w(r,0)=G(r)$; see, for example,~\cite[Section~3,
equation~(11)]{Tao}. The cited formula uses frequency $\eta$ with exponential $e^{-2\pi i r\eta}$;
the change $\xi=2\pi\eta$ gives the factor $1/(2\pi)$ in \eqref{eq:heat}.
Since $G$ is real and even, so is $w$.
Moreover
\[
 \sup_r|w(r,t)-G(r)|\le\frac1{2\pi}\int_\R
 |e^{-\kappa t|\xi|^p}-1|\,|\widehat G(\xi)|\,d\xi\longrightarrow0
\]
by dominated convergence.

For each frequency, the absolute value of the symmetric exponential
difference is $2(1-\cos(z\xi))$. Its jump integral equals
$(\kappa/C_s)|\xi|^p$, with value zero at $\xi=0$.
The double absolute integral is therefore bounded by
$(\kappa/C_s)\int_\R|\xi|^p|\widehat G(\xi)|\,d\xi<\infty$.
Fubini's theorem, followed by $z|\xi|$ substitution, gives
\[
 \Lop w(r,t)=-\frac1{2\pi}\int_\R
 e^{ir\xi-\kappa t|\xi|^p}\kappa|\xi|^p\widehat G(\xi)\,d\xi=w_t(r,t).
\]
This also proves $w_t(r,0)=b(r)$. The Fourier formula is consistent with the
one-dimensional fractional heat representation in~\cite[Section~5]{DEJV};
no monotonicity hypothesis from its Lemma~5.1 is used here.

Evenness and the preceding radial Hessian argument give bounded continuous
spatial derivatives through order two for $V$, including at zero. The
Fourier majorants make these derivatives and $V_t$ jointly continuous.
For $x\ne0$ its radial directional integral is
\[
 D_{e_r}^sV(x,t)=\Lop w(|x|,t)=V_t(x,t).
\]
At a nonzero gradient this is the selected direction, up to sign. At a zero
gradient the radial value lies between the infimum and supremum over all
directions. At the origin every directional value equals $V_t$. This proves
both pointwise inequalities, even at noncentral critical radii.
\end{proof}

\begin{proof}[Proof of Theorem~\ref{thm:main}]
Use the $U$ and $V$ constructed above, restricted to $[0,T]$. They have the
claimed regularity and boundedness on this finite cylinder.
We first pass from the pointwise inequalities to every actual viscosity test.
At an upper contact with either constructed function, a local test has the same
spatial gradient and time derivative at contact. On that time slice its patch
is everywhere greater than or equal to the tested function and agrees at the
center. Hence its directional integral is greater than or equal to the
function's integral in every direction. The selected-direction inequality,
or its supremum at zero gradient, follows. A lower patch reverses every
directional inequality, giving the supersolution condition with the infimum.
The integrals converge absolutely by local $C^2$ control near contact and a
bounded tail. Lemmas~\ref{lem:sub} and~\ref{lem:heat} therefore give precisely
the subsolution and supersolution properties required in the theorem.

Their common initial profile is $F\in C_c^\infty$, hence bounded and uniformly
continuous. The trace for $V$ is uniform by Lemma~\ref{lem:heat}; for $U$ use
$\sup_x|U(x,t)-F(x)|\le t\norm Q+Mt^2$.
Finally set $x_*=2e_1$, where $e_1$ is the first coordinate vector, and
$\delta=a-b(2)>0$. The Fourier formula gives the finite bound
\[
 |w_{tt}(r,t)|\le K_2:=\frac1{2\pi}\int_\R
 \kappa^2|\xi|^{2p}|\widehat G(\xi)|\,d\xi.
\]
Since $G(2)=0$ and $q(2)=a$, Taylor's formula in time yields
\[
 U(x_*,t)-V(x_*,t)\ge\delta t-(M+K_2/2)t^2.
\]
Choose
\[
 t_*=\min\{T/2,\delta/[2(M+K_2/2)]\},\qquad 0<t_*<T.
\]
Then $U(x_*,t_*)-V(x_*,t_*)\ge\delta t_*/2>0$, as required.
\end{proof}

\paragraph{AI use.}
AI tools assisted with the mathematical construction, proof development,
source checking, and manuscript preparation. The manuscript is submitted for
independent mathematical review.

\begin{thebibliography}{9}
\bibitem{BCF}
C.~Bjorland, L.~Caffarelli and A.~Figalli,
\emph{Non-Local Tug-of-War and the Infinity Fractional Laplacian},
arXiv:1011.1966v2 (2011),
\url{https://arxiv.org/abs/1011.1966v2}.
\bibitem{DEJV}
F.~del Teso, J.~Endal, E.~R.~Jakobsen and J.~L.~V\'azquez,
\emph{Evolution Driven by the Infinity Fractional Laplacian},
arXiv:2210.06414v1 (2022),
\url{https://arxiv.org/abs/2210.06414v1}.
\bibitem{Tao}
T.~Tao, \emph{245C, Notes 2: The Fourier transform},
6 April 2009, Section~3, equation~(11),
\url{https://terrytao.wordpress.com/2009/04/06/the-fourier-transform/}.
\end{thebibliography}
\end{document}
